\section{The algebra and its commutant}
\label{sec:algebra}

\subsection{Finite weighted power systems}
\label{sec:systems}

The finite weighted dynamics that the algebra linearizes is the following.

\begin{definition}
\label{def:weighted-power-system}
Fix a finite label set \(E\).  A \emph{finite weighted power system} is a
tuple
\[
 \mathscr X=(X,\mu,x_0,(\tau_e)_{e\in E}),
\]
where \(X\) is a finite set, \(\mu:X\to\Q_{>0}\) has total mass one,
\(x_0\in X\), and the maps \(\tau_e:X\to X\) commute and fix \(x_0\).
An isomorphism \(a:\mathscr X\to\mathscr Y\) is a bijection satisfying
\[
 a(x_0)=y_0,\qquad \nu(a(x))=\mu(x),\qquad
 a\tau_e=\upsilon_ea\quad(e\in E).
\]
\end{definition}

\begin{definition}
\label{def:multiplicative-labels}
Call the system \emph{multiplicatively labeled} if \(E\) is a set of positive
integers, \(1\in E\), and
\begin{equation}
 \tau_m\tau_n=\tau_{mn}\ \text{whenever }m,n,mn\in E,
 \qquad
 \tau_1=\mathrm{id}.
 \label{eq:label-multiplicativity}
\end{equation}
This hypothesis is \emph{not} part of \cref{def:weighted-power-system}; it is
invoked explicitly where it is used, namely in
\cref{prop:minimal-prime-generators} and throughout
\cref{sec:goodlabels}.  The structural results of the present section do
not require it.
\end{definition}

Write \(V_X=\operatorname{Map}(X,\Q)\), let
\(\mathcal B_X=(\delta_x)_{x\in X}\) be its coordinate basis, and put
\[
 \langle f,h\rangle_\mu=\sum_{x\in X}\mu(x)f(x)h(x).
\]
The pullback \(P_e f=f\circ\tau_e\) has coordinate matrix
\begin{equation}
 (P_e)_{xy}=\mathbf1_{\{y=\tau_e(x)\}},
 \qquad
 P_e^*=D_\mu^{-1}P_e^{\mathsf T}D_\mu,
 \quad D_\mu=\operatorname{diag}(\mu(x)).
 \label{eq:system-adjoint}
\end{equation}
Let \(\mathbf c_\mu=(c_x)_{x\in X}\) be the primitive positive integral
vector proportional to \(\mu\), and regard both
\(\mathbf c_\mu=\sum_xc_x\delta_x\) and
\(\eta_0=\delta_{x_0}\) as marked vectors.

\begin{definition}
\label{def:fully-marked-object}
The fully marked linear object attached to \(\mathscr X\) is
\[
 \mathbf P(\mathscr X)=
 \bigl(
 V_X,\mathcal B_X,\langle\ ,\ \rangle_\mu,\mathbf1,\eta_0,
 \mathbf c_\mu,(P_e)_{e\in E},
 \PE(\mathscr X)
 \bigr),
\]
where
\[
 \PE(\mathscr X)
 =\Alg_\Q\{I,P_e,P_e^*:e\in E\}\subseteq\End_\Q(V_X).
\]
\end{definition}

The forward subalgebra
\(\SfwdE(\mathscr X)=\Alg_\Q\{I,P_e:e\in E\}\) is commutative, since the
\(\tau_e\) commute; \(\PE(\mathscr X)\) is its star closure.
If \(a:\mathscr X\to\mathscr Y\) is an isomorphism, then
\((U_af)(y)=f(a^{-1}y)\) carries every part of the marking to the
corresponding part, is an isometry, and intertwines \(P_e^X\) with
\(P_e^Y\) and their adjoints.  Conversely an isomorphism of fully marked
objects permutes their coordinate bases and therefore comes from a unique
system isomorphism.  The fully marked object thus reconstructs
\((X,\mu,x_0,(\tau_e))\) up to unique simultaneous relabeling.

More general factor maps do not act on the star closure.  Suppose
\(q:(X,\mu,\tau_e)\to(Y,\nu,\upsilon_e)\) is a surjection with
\(q_*\mu=\nu\) and \(q\tau_e=\upsilon_eq\).  Pullback
\(U_q:V_Y\to V_X\), \(U_qf=f\circ q\), is an isometry and
\begin{equation}
 P_e^XU_q=U_qP_e^Y.
 \label{eq:factor-forward-intertwining}
\end{equation}
It follows only that
\(U_q^*(P_e^X)^*=(P_e^Y)^*U_q^*\).  The stronger relation
\((P_e^X)^*U_q=U_q(P_e^Y)^*\), needed to intertwine the generated
adjoint-closed algebras, holds if and only if
\(\operatorname{im}U_q\) is reducing for \(P_e^X\).  Equivalently, for
every \(x\in X\) and \(z\in Y\),
\begin{equation}
 \frac1{\mu(x)}
 \sum_{\substack{\tau_e y=x\\q(y)=z}}\mu(y)
 =
 \begin{cases}
 \displaystyle\frac{\nu(z)}{\nu(q(x))},
   &\upsilon_ez=q(x),\\[6pt]
 0,&\upsilon_ez\ne q(x).
 \end{cases}
 \label{eq:reverse-fibre-condition}
\end{equation}
General measure-preserving factor maps need not satisfy this
reverse-fibre condition.

\subsubsection{The system supplied by a finite group}

Let
\(\operatorname{Conj}(G)=\{C_1,\ldots,C_k\}\), with
\(C_1=\{1\}\), and set
\begin{equation}
 V_G=\operatorname{Map}(\operatorname{Conj}(G),\Q),\qquad
 \langle f,h\rangle_G
 =\frac1{|G|}\sum_{i=1}^k |C_i|f(C_i)h(C_i).
 \label{eq:coordinate-class-form}
\end{equation}
For \(n\geq1\), write \(\sigma_n(i)=j\) when \(g_i^n\in C_j\).
Then
\begin{equation}
 (P_n)_{ij}=\mathbf1_{\{j=\sigma_n(i)\}},
 \qquad
 P_n^*=D_G^{-1}P_n^{\mathsf T}D_G,\quad
 D_G=\operatorname{diag}\left(\frac{|C_i|}{|G|}\right).
 \label{eq:power-matrix}
\end{equation}
Thus \(G\) supplies the finite weighted power system
\[
 \mathscr C_E(G)=
 \left(\operatorname{Conj}(G),\frac{|C|}{|G|},\{1\},
 ([x]\mapsto[x^n])_{n\in E}\right).
\]
Its primitive weight vector is the positive class-size vector
\((|C_i|)_i\).  We define
\begin{equation}
 \mathbf P(G,E)
 =\mathbf P(\mathscr C_E(G)),\qquad
 \PE(G)
 =\PE(\mathscr C_E(G)).
 \label{eq:ap-definition-structural}
\end{equation}
For \(N=\exp(G)\), the canonical label set is
\(E=\pi(N)\), and we omit \(E\) from the notation.

A group isomorphism induces an isomorphism of these fully marked objects.
More generally, simultaneous relabeling of classes preserving class sizes,
the identity class, and all selected power maps is isomorphism at
the source-system level.

\begin{proposition}
\label{prop:quotient-factor}
If \(q:G\twoheadrightarrow Q\) is a quotient, then
\[
 \bar q:\operatorname{Conj}(G)\longrightarrow\operatorname{Conj}(Q),
 \qquad [g]\longmapsto[q(g)],
\]
is a surjective measure-preserving factor map for every power map.
Consequently pullback \(U_{\bar q}:V_Q\to V_G\) is an isometry and
\(\Psi_n^GU_{\bar q}=U_{\bar q}\Psi_n^Q\).
It intertwines the full algebras if and only if its image
is reducing for every \(\Psi_n^G\).
\end{proposition}

\begin{proof}
The map on classes is well defined, surjective, and commutes with powers.
For a class \(D\subseteq Q\), its full inverse image has
\(|\ker q||D|\) elements.  Hence the sum of the normalized \(G\)-class
weights above \(D\) is
\[
 \frac{|\ker q||D|}{|G|}=\frac{|D|}{|Q|}.
\]
The remaining assertions are \eqref{eq:factor-forward-intertwining} and its
reverse-fibre criterion.
\end{proof}

A proper subgroup inclusion does not preserve normalized class measure
(the identity class has mass \(1/|H|\) in \(H\) and \(1/|G|\) in \(G\)),
so no proper inclusion is a morphism of weighted power systems.

\subsection{Transfer operators}
\label{sec:transfer}

Let \(\mathscr X=(X,\mu,x_0,(\tau_e)_{e\in E})\) be a finite weighted
power system.  Write
\[
 P_ef=f\circ\tau_e,\qquad L_e=P_e^*,
\]
and let \(M_a\) denote multiplication by \(a\in V_X\).  The reverse-fibre
density is
\begin{equation}
 r_e(y)=\frac{1}{\mu(y)}
 \sum_{\tau_ex=y}\mu(x).
 \label{eq:reverse-fibre-density}
\end{equation}

\begin{theorem}
\label{thm:transfer-calculus}
For \(a,b\in V_X\) and \(e,f\in E\),
\begin{align}
 P_eM_a&=M_{P_ea}P_e,                         \label{eq:covariance}\\
 L_e\bigl((P_ea)b\bigr)&=aL_eb,               \label{eq:frobenius-transfer}\\
 L_eM_bP_e&=M_{L_eb},                         \label{eq:transfer-conjugation}\\
 L_eP_e&=M_{r_e}.                             \label{eq:gram-density}
\end{align}
If \(\tau_{ef}=\tau_e\tau_f=\tau_f\tau_e\), then
\begin{equation}
 L_eL_f=L_fL_e=L_{ef},\qquad
 r_{ef}=L_er_f=L_fr_e.
 \label{eq:density-cocycle}
\end{equation}
\end{theorem}

\begin{proof}
The first identity follows from
\[
 (P_e(ah))(x)=a(\tau_ex)h(\tau_ex).
\]
The adjoint is
\begin{equation}
 (L_eb)(y)=\frac{1}{\mu(y)}
 \sum_{\tau_ex=y}\mu(x)b(x).
 \label{eq:explicit-transfer}
\end{equation}
On the fibre \(\tau_ex=y\), the function \(P_ea\) has value \(a(y)\).
This gives \eqref{eq:frobenius-transfer} and
\eqref{eq:transfer-conjugation}.  Taking \(b=1\) gives
\eqref{eq:gram-density}.  The last two identities follow by taking
adjoints in \(P_eP_f=P_fP_e=P_{ef}\) and applying the result to \(1\).
\end{proof}

\begin{theorem}
\label{thm:normality-commutativity}
For a label \(e\), the following conditions are equivalent:
\begin{enumerate}
\item \(P_e\) is normal;
\item \(P_e\) is unitary;
\item \(\tau_e\) is a measure-preserving permutation;
\item \(r_e=1\) and \(\tau_e\) is injective.
\end{enumerate}
Consequently
\begin{equation}
 \PE(\mathscr X)\text{ is commutative}
 \quad\Longleftrightarrow\quad
 \text{every \(\tau_e\) is a measure-preserving permutation.}
 \label{eq:commutative-dichotomy}
\end{equation}
\end{theorem}

\begin{proof}
The matrix \(L_eP_e=M_{r_e}\) is diagonal, while
\begin{equation}
 (P_eL_e)_{xz}=
 \begin{cases}
 \mu(z)/\mu(\tau_ex),&\tau_ez=\tau_ex,\\
 0,&\tau_ez\neq\tau_ex.
 \end{cases}
 \label{eq:fibre-gram}
\end{equation}
If \(P_e\) is normal, \eqref{eq:fibre-gram} is diagonal.  Every fibre of
\(\tau_e\) then has at most one point, so \(\tau_e\) is a permutation.
For a permutation,
\[
 r_e(y)=\frac{\mu(\tau_e^{-1}y)}{\mu(y)}.
\]
Normality gives \(r_e(x)=r_e(\tau_ex)\).  This value is constant on each
cycle, and its product around the cycle is \(1\); hence \(r_e=1\).
Thus \(\tau_e\) preserves \(\mu\) and \(P_e\) is unitary.  The remaining
implications are immediate.  Commuting measure-preserving permutations
have commuting inverse pullbacks, which proves
\eqref{eq:commutative-dichotomy}.
\end{proof}

\begin{corollary}
\label{cor:group-normal-labels}
Let \(G\) be finite and \(N=\exp(G)\).  Then
\begin{equation}
 \Psi_n\text{ is normal}
 \quad\Longleftrightarrow\quad
 \Psi_n\text{ is unitary}
 \quad\Longleftrightarrow\quad
 (n,N)=1.
 \label{eq:good-normal-unitary}
\end{equation}
The good labels generate a commutative semisimple algebra.  Every prime
\(p\mid N\) gives a nonnormal generator; hence
\(\Pram(G)\) is noncommutative for \(G\neq1\).
\end{corollary}

\begin{proof}
If \((n,N)=1\), choose \(m\) with \(mn\equiv1\pmod N\).  The \(n\)-th and
\(m\)-th power maps are inverse permutations of the conjugacy classes and
preserve their sizes.  If \((n,N)>1\), choose
\(p\mid(n,N)\) and an element \(x\) of order \(p\).  Then \(x^n=1\), so
the power map on conjugacy classes is not injective.  Apply
\cref{thm:normality-commutativity}.
\end{proof}

For a finite group, \eqref{eq:reverse-fibre-density} becomes
\begin{equation}
 r_n(C)=
 \frac{\#\{x\in G:x^n\in C\}}{|C|},
 \qquad
 r_n(\{1\})=\#\{x\in G:x^n=1\}.
 \label{eq:group-ramification-density}
\end{equation}
Thus \(\Psi_n^*\Psi_n\) is the diagonal operator of class-resolved
\(n\)-th root counts.

\begin{definition}
\label{def:ramification-algebra}
Let \(\mathcal R_E(\mathscr X)\) be the smallest unital
\(\Q\)-subalgebra of \(V_X\) stable under every \(L_e\), and put
\[
 \mathcal D_E(\mathscr X)
 =\{M_a:a\in\mathcal R_E(\mathscr X)\}.
\]
\end{definition}

\begin{proposition}
\label{prop:ramification-diagonal-inclusion}
For every finite weighted power system,
\[
 \mathcal D_E(\mathscr X)\subseteq\PE(\mathscr X).
\]
\end{proposition}

\begin{proof}
The identity multiplier belongs to \(\PE(\mathscr X)\).  If \(M_a\)
belongs to it, then
\[
 M_{L_ea}=L_eM_aP_e\in\PE(\mathscr X)
\]
by \eqref{eq:transfer-conjugation}.  The multipliers are closed under sums
and products.
\end{proof}

\begin{theorem}
\label{thm:diagonal-separation}
Let \(\Gamma_E\) be the undirected graph on \(X\) with edges
\(\{x,\tau_ex\}\).  If \(\mathcal R_E(\mathscr X)\) separates the points
of \(X\), then
\begin{equation}
 \PE(\mathscr X)
 \cong\bigoplus_{C\in\pi_0(\Gamma_E)}
       \End_\Q(\Q^C).
 \label{eq:component-matrix-decomposition}
\end{equation}
If \(\Gamma_E\) is connected, then
\(\PE(\mathscr X)=\End_\Q(V_X)\).
\end{theorem}

\begin{proof}
For \(x\in X\) and \(y\neq x\), choose \(a_y\in\mathcal R_E(\mathscr X)\)
with \(a_y(x)\neq a_y(y)\).  Then
\[
 q_x=\prod_{y\neq x}
 \frac{a_y-a_y(y)}{a_y(x)-a_y(y)}
 =\mathbf1_{\{x\}}.
\]
Hence all coordinate idempotents \(E_{xx}=M_{\mathbf1_{\{x\}}}\) belong
to \(\PE(\mathscr X)\).  If \(\tau_ex=y\), then
\[
 E_{xx}P_eE_{yy}=E_{xy},\qquad
 E_{yy}L_eE_{xx}=\frac{\mu(x)}{\mu(y)}E_{yx}.
\]
Products along paths in \(\Gamma_E\) give every matrix unit within each
connected component.  Every generator preserves the direct sum over the
components.
\end{proof}

\subsection{Semisimplicity and the rational bicommutant}
\label{sec:structural-theory}

\begin{proposition}
\label{prop:minimal-prime-generators}
Let \(N=\exp(G)\).  Each of the following sets generates the same unital
adjoint-closed rational algebra:
\begin{equation}
 \{\Psi_p:p\mid N,\ p\text{ prime}\},\qquad
 \{\Psi_{p^a}:p^a\mid N\},\qquad
 \{\Psi_d:d\mid N\}.
 \label{eq:equivalent-exponent-sets}
\end{equation}
Here the adjoints of the displayed operators are included in every case.
This common algebra is \(\Pram(G)\).
\end{proposition}

\begin{proof}
The identities
\[
 \Psi_m\Psi_n=\Psi_{mn},\qquad \Psi_1=I
\]
show that \(\Psi_{p^a}=\Psi_p^a\) and, for
\(d=\prod_p p^{a_p}\), that \(\Psi_d=\prod_p\Psi_p^{a_p}\).
Thus all divisor maps lie in the algebra generated by the prime maps.  The
reverse inclusion is immediate.  Since an adjoint-closed algebra contains
the adjoint of every product it contains, adjoining the listed adjoints
does not alter this conclusion.
\end{proof}

\begin{example}
\label{ex:all-positive-exponents}
The algebra generated by \(\Psi_n\) for every \(n\geq1\) need not equal
\(\Pram(G)\).  For example, if \(G=C_5\), the canonical
algebra uses only \(\Psi_5\), whereas \(\Psi_2\) cyclically permutes the
four nonidentity class coordinates.  The prime-exponent calculation in
\cref{thm:prime-exponent-structural} shows that the canonical algebra acts
scalarly on a three-dimensional complementary summand, while \(\Psi_2\)
does not.  Hence the all-exponents algebra is strictly larger in this
example.  Proposition~\ref{prop:minimal-prime-generators} concerns the
power maps indexed by divisors of \(\exp(G)\), not all residue classes
modulo the exponent.
\end{example}

For a subalgebra \(A\subseteq\End_F(V)\), write \(A'\) for its commutant
in \(\End_F(V)\).

\begin{lemma}
\label{lem:commutant-base-change}
Let \(F\subseteq K\) be a field extension, let \(V\) be a
finite-dimensional \(F\)-vector space, and let
\(A\subseteq\End_F(V)\) be a finite-dimensional subalgebra.  Under the
natural identification
\(\End_F(V)\otimes_FK\cong\End_K(V\otimes_FK)\), one has
\begin{equation}
 A'\otimes_FK=(A\otimes_FK)'.
 \label{eq:commutant-base-change}
\end{equation}
\end{lemma}

\begin{proof}
Choose an \(F\)-basis \(a_1,\ldots,a_r\) of \(A\).  Its commutant is the
kernel of the \(F\)-linear map
\[
 T:\End_F(V)\longrightarrow\End_F(V)^{\oplus r},
 \qquad X\longmapsto([X,a_1],\ldots,[X,a_r]).
\]
Tensoring the exact sequence
\(0\to\ker T\to\End_F(V)\xrightarrow{T}\End_F(V)^{\oplus r}\)
with the field \(K\) preserves exactness.  The resulting kernel is
the commutant of \(A\otimes_FK\).
\end{proof}

\begin{theorem}
\label{thm:rational-bicommutant-structural}
Let \(X\) be a finite set, let \(\mu:X\to\Q_{>0}\), and let
\((T_e)_{e\in E}\) be any finite family of \(\Q\)-linear operators on
\(V_X\), with \(T_e^*=D_\mu^{-1}T_e^{\mathsf T}D_\mu\) the adjoint for
\(\langle\ ,\ \rangle_\mu\).  Then
\[
 A=\Alg_\Q\{I,T_e,T_e^*:e\in E\}\subseteq\End_\Q(V_X)
\]
is a finite-dimensional semisimple \(\Q\)-algebra, and in its defining
faithful representation on \(V_X\),
\begin{equation}
 A''=A.
 \label{eq:rational-double-commutant-structural}
\end{equation}
Over \(\C\), there are pairwise inequivalent simple modules
\(W_1,\ldots,W_t\), with \(d_j=\dim_\C W_j\), and multiplicity spaces
\(M_j\), with \(m_j=\dim_\C M_j\), such that
\begin{equation}
 V_X\otimes_\Q\C\cong\bigoplus_{j=1}^t W_j\otimes_\C M_j,
 \qquad
 A\otimes_\Q\C\cong
 \bigoplus_{j=1}^t\End_\C(W_j)\otimes I_{M_j}.
 \label{eq:complex-isotypic-structural}
\end{equation}
In particular,
\begin{equation}
 \dim_\C(A\otimes\C)=\sum_jd_j^2,
 \qquad
 \dim_\C(A\otimes\C)'=\sum_jm_j^2.
 \label{eq:dimension-identities-structural}
\end{equation}
In particular this applies to \(A=\PE(\mathscr X)\) for any finite weighted
power system, and to \(A=\PE(G)\) for any finite group and any exponent set.
\end{theorem}

\Cref{thm:rational-bicommutant-structural} uses none of the structure of a
finite weighted power system beyond positivity of \(\mu\): the \(T_e\) need
not commute, need not be pullbacks along maps, and no distinguished point is
required.  The star closure alone gives semisimplicity and
bicommutancy.  Deterministic commuting maps are what give the balanced-kernel
description of the commutant (\cref{thm:balanced-commutant}), and
multiplicative labels are what give the good-label arithmetic of
\cref{sec:goodlabels}.

\begin{proof}
The algebra is finite-dimensional because it is a subalgebra of
\(\End_\Q(V_X)\).  Extend scalars and the form \(\langle\ ,\ \rangle_\mu\) to
obtain a Hermitian form on \(V_X\otimes_\Q\C\), positive definite because
\(\mu>0\) pointwise.  The complex
algebra \(A_\C=A\otimes_\Q\C\) is closed under its Hilbert-space adjoint.
Its Jacobson radical \(J\) is preserved by the involution because the
involution is an anti-automorphism and the radical is invariant under
anti-automorphisms.  If \(x\in J\), then \(x^*x\in J\).  The ideal \(J\)
is nilpotent, while \(x^*x\) is positive semidefinite, so \(x^*x=0\) and
therefore \(x=0\).  Hence \(A_\C\), and consequently \(A\), is
semisimple.

The finite-dimensional structure theorem gives the
isotypic form \eqref{eq:complex-isotypic-structural}~\cite{Wedderburn}.
On each summand,
its commutant is
\(I_{W_j}\otimes\End_\C(M_j)\), whose commutant is
\(\End_\C(W_j)\otimes I_{M_j}\).  Thus \(A_\C''=A_\C\).  Applying
\cref{lem:commutant-base-change} twice gives
\[
 A''\otimes_\Q\C=A_\C''=A_\C=A\otimes_\Q\C.
\]
Since \(A\subseteq A''\) and scalar extension preserves dimensions, this
forces \(A=A''\).  The dimension identities follow from the two displayed
block actions.  The final sentence is the case \(T_e=P_e\) of
\eqref{eq:system-adjoint}.
\end{proof}

\begin{proposition}
\label{prop:field-extension}
For every field extension \(K/\Q\),
\[
 \PE(G)\otimes_\Q K
 =\Alg_K\{I,\Psi_n,\Psi_n^*:n\in E\}
 \subseteq\End_K(V_G\otimes_\Q K),
\]
and
\[
 \PE(G)'\otimes_\Q K
 =(\PE(G)\otimes_\Q K)'.
\]
Thus extension of scalars preserves the represented algebra and its
commutant.  The center fields, Galois orbits, and Schur indices give the
descent data that recover the rational form from its complex blocks.
\end{proposition}

\begin{proof}
The first identity follows because the same rational generators span both
sides after scalar extension.  The second is
\cref{lem:commutant-base-change}.  The final statement is the usual
distinction between a semisimple rational algebra and its split complex
form.
\end{proof}

Direct products exhibit a natural tensor relation, but the power
maps on the two factors are synchronized by the same exponent.

\begin{proposition}
\label{prop:direct-product-containment}
For finite groups \(G,H\) and any exponent set \(E\), the canonical
identification
\[
 V_{G\times H}\cong V_G\otimes_\Q V_H
\]
satisfies
\begin{equation}
 \Psi_n^{G\times H}=\Psi_n^G\otimes\Psi_n^H,
 \qquad
 (\Psi_n^{G\times H})^*=(\Psi_n^G)^*\otimes(\Psi_n^H)^*.
 \label{eq:direct-product-generators}
\end{equation}
Consequently,
\begin{equation}
 \PE(G\times H)
 \subseteq
 \PE(G)\otimes_\Q
 \PE(H)
 \label{eq:direct-product-containment}
\end{equation}
as represented algebras.  For canonical exponent sets, take
\(E=\pi(\exp(G\times H))\).  If \(G\) and \(H\) have the same prime
divisors in their exponents, the right-hand side is
\(\Pram(G)\otimes
\Pram(H)\).
\end{proposition}

\begin{proof}
The conjugacy classes of \(G\times H\) are the products \(C_i\times D_j\),
their sizes multiply, and
\((g,h)^n=(g^n,h^n)\).  Hence the power matrix and class-measure matrix are
\[
 P_n^{G\times H}=P_n^G\otimes P_n^H,
 \qquad
 D_{G\times H}=D_G\otimes D_H.
\]
Formula \eqref{eq:system-adjoint} gives the second identity in
\eqref{eq:direct-product-generators}.  Every generator on the left of
\eqref{eq:direct-product-containment} therefore belongs to the tensor
product algebra on the right.
\end{proof}

Containment in \eqref{eq:direct-product-containment} can be strict.  The
left-hand algebra is generated by the synchronized tensors
\(\Psi_n^G\otimes\Psi_n^H\) and their adjoints.  Tensor-product equality
holds if and only if these synchronized generators recover the separate
operators \(\Psi_n^G\otimes I\) and \(I\otimes\Psi_n^H\); strict examples
occur when they do not.

\subsection{Balanced kernels}
\label{sec:reconstruction}

The commutant has an intrinsic description directly on the weighted colored
digraph.  Write an operator \(S\in\End_\Q(V_X)\) in coordinate form
\[
 (Sf)(x)=\sum_{y\in X}s(x,y)f(y).
\]

\begin{definition}
\label{def:balanced-kernels}
Let \(\operatorname{Bal}(\mathscr X)\) be the rational vector space of
kernels \(s:X\times X\to\Q\) satisfying, for every \(e\in E\) and
\(x,z\in X\),
\begin{align}
 s(\tau_e x,z)
 &=\sum_{\tau_e y=z}s(x,y),
 \label{eq:forward-balance}\\
 \frac{\mu(z)}{\mu(\tau_e z)}s(x,\tau_e z)
 &=\frac1{\mu(x)}
   \sum_{\tau_e y=x}\mu(y)s(y,z).
 \label{eq:reverse-balance}
\end{align}
Composition is ordinary matrix multiplication.
\end{definition}

These kernels balance the colored forward fibres and their weighted reverse
fibres.

\begin{theorem}
\label{thm:balanced-commutant}
For every finite weighted power system,
\begin{equation}
 \PE(\mathscr X)'=\operatorname{Bal}(\mathscr X).
 \label{eq:balanced-commutant}
\end{equation}
Consequently
\begin{equation}
 \PE(\mathscr X)
 =\operatorname{Bal}(\mathscr X)'
 \label{eq:balanced-bicommutant}
\end{equation}
in its defining representation.  In particular the algebra is determined
by the fractional symmetries of the weighted colored power graph.
\end{theorem}

\begin{proof}
Matrix multiplication gives
\[
 (SP_e)_{xz}=\sum_{\tau_e y=z}s(x,y),\qquad
 (P_eS)_{xz}=s(\tau_ex,z),
\]
so \(SP_e=P_eS\) is \eqref{eq:forward-balance}.  From
\((P_e^*)_{yz}=\mu(z)\mu(y)^{-1}\mathbf1_{\{y=\tau_ez\}}\), one obtains
\[
 (SP_e^*)_{xz}=\frac{\mu(z)}{\mu(\tau_ez)}s(x,\tau_ez),\qquad
 (P_e^*S)_{xz}=\frac1{\mu(x)}
 \sum_{\tau_ey=x}\mu(y)s(y,z).
\]
Thus \(SP_e^*=P_e^*S\) is \eqref{eq:reverse-balance}.  Commuting with all
generators is equivalent to commuting with their generated algebra, proving
\eqref{eq:balanced-commutant}.  Equation
\eqref{eq:balanced-bicommutant} is the rational double-commutant theorem.
\end{proof}

Every source automorphism gives a permutation kernel in
\(\operatorname{Bal}(\mathscr X)\).  Automorphisms alone usually do not
describe the commutant: arbitrary rational balanced
kernels are allowed.  For a group of prime exponent with \(k\) classes, the
decomposition in \cref{thm:prime-exponent-structural} gives the explicit
formula
\[
 \Pram(G)'
 \cong \Q\oplus M_{k-2}(\Q),\qquad
 \dim_\Q\Pram(G)'=1+(k-2)^2.
\]
This is a family-level commutant calculation rather than a matrix nullspace
computation.

Between the fully marked object and the abstract algebra there are two
further levels of forgetting: the operator-marked object
\((V,\langle\ ,\ \rangle,\mathbf1,\eta_0,(P_e))\), which retains the
labeled generators but not the coordinate basis, and the unmarked
represented algebra \((A\subseteq\End_\Q(V),V,\langle\ ,\ \rangle)\).
Each arrow in this tower is strict over \(\Q\): the last by the
prime-exponent family, where \(C_3\) and \(C_5\) share the abstract algebra
\(M_2(\Q)\oplus\Q\) but have scalar factors of different multiplicity; the
others by explicit small systems, which we omit.

\subsection{Saturation and two degeneracy criteria}
\label{sec:saturation}

The balanced-kernel theorem converts three questions about the algebra
into questions about balanced kernels on the power graph: when enlarging the
label set changes nothing, when the algebra degenerates to everything, and
when the commutant is commutative.  All three are consequences of
\cref{thm:balanced-commutant} together with the rational double commutant.

\begin{theorem}
\label{thm:saturation}
Let \(E\subseteq F\) be finite label sets for \(\mathscr X\).  Then
\(\PE(\mathscr X)\subseteq\PF(\mathscr X)\), and
\begin{equation}
 \PE(\mathscr X)=\PF(\mathscr X)
 \iff
 \text{every }P_f,\ f\in F\setminus E,\ \text{commutes with every}
 \ s\in\operatorname{Bal}_E(\mathscr X),
 \label{eq:saturation}
\end{equation}
where \(\operatorname{Bal}_E\) is the balanced-kernel algebra of
\cref{def:balanced-kernels} formed with the label set \(E\).  We call \(E\)
\emph{saturated} in \(F\) when this holds.
\end{theorem}

\begin{proof}
The containment is clear.  For the criterion, \(\PF\) is generated
over \(\PE\) by the operators \(P_f\) and \(P_f^*\) with \(f\in
F\setminus E\), so equality holds if and only if every such \(P_f\) lies in
\(\PE\); the adjoints then follow because \(\PE\) is adjoint closed.  By
\cref{thm:rational-bicommutant-structural} and
\cref{thm:balanced-commutant},
\[
 \PE=(\PE)''=\operatorname{Bal}_E(\mathscr X)',
\]
and membership in the commutant of \(\operatorname{Bal}_E\) is the
stated commutation condition.
\end{proof}

\begin{corollary}
\label{cor:full-matrix}
\(\PE(\mathscr X)=\End_\Q(V_X)\) if and only if
\(\operatorname{Bal}_E(\mathscr X)=\Q\), that is, if and only if the only
balanced kernels are the scalar multiples of the identity kernel.
\end{corollary}

\begin{proof}
By \cref{thm:balanced-commutant} and the double commutant,
\(\PE=\operatorname{Bal}_E'\), and a subalgebra of \(\End_\Q(V_X)\) equals
\(\End_\Q(V_X)\) if and only if its commutant is the scalars.
\end{proof}

\begin{corollary}
\label{cor:multiplicity-free}
The defining module \(V_X\) is multiplicity free over
\(\PE(\mathscr X)\otimes\C\) if and only if \(\operatorname{Bal}_E(\mathscr X)\) is
commutative, and in that case
\(\dim_\Q Z(\PE)=\dim_\Q\operatorname{Bal}_E(\mathscr X)\).
\end{corollary}

\begin{proof}
With notation as in \eqref{eq:complex-isotypic-structural}, the commutant
after scalar extension is \(\bigoplus_j\End_\C(M_j)\), which is commutative
if and only if every \(m_j=1\).  Commutativity is preserved and reflected by
scalar extension, so it may be tested over \(\Q\).  When all \(m_j=1\) the
commutant and the center both have dimension \(t\).
\end{proof}

Throughout, ``\(V_G\) is multiplicity free'' means the condition of
\cref{cor:multiplicity-free}: the commutant \(\Pram(G)'\) is commutative.

First, \cref{thm:saturation} explains \cref{ex:all-positive-exponents}: for
\(G=C_5\) the label set \(E=\{5\}\) is not saturated in \(\{2,5\}\), because
\(\operatorname{Bal}_{\{5\}}\) contains the full matrix algebra of the
three-dimensional complementary summand, and the four-cycle \(\Psi_2\) does
not commute with all of it.  Hence
\(\Pram(C_5)\subsetneq\PAd(C_5)\), with dimensions \(5\) and \(7\).
Second, by \cref{cor:full-matrix} the equality \(\dim A=k(G)^2\) holds when
the weighted colored power graph admits no balanced kernel beyond
the scalars.  The rows of \cref{thm:comparison-panel} with
\(\dim A=k(G)^2\) therefore record rigidity of that graph.
